\scanpage{129}
\mspage{24 bis}
\deleted{
\deleted{
$G_b$, où $b\in\Gamma(G,O_G)$ est la section définie dans 2.7. (NB on \unsure{décomposera} $S$ en le préschéma somme des préschémas $S_{n,\rho}$, sur chacun desquels $G$ est de dimension relative $n$ et de rang \deleted{réductif} nilpotent $\rho$, et on construit $b$ séparément sur chacun des $G\vert S_{n,\rho}$). D'autre part, comme $\unsure{\Psi|\psi}^{-1}(U)$ et $U$ sont lisses donc plats localement de présentation finie sur $S$, pour vérifier que $\psi$ induit un isomorphisme $\psi^{-1}(U)\longrightarrow U$, il suffit de le vérifier fibre par fibre, \inserted{(SGA I 5.7.)} ce qui nous ramène au cas où $S$ est le spectre d'un corps, que l'on peut supposer algébriquement clos. Il existe alors un sous-groupe de Cartan $C$ de $G$, et si $N$ est le normalisateur, on sait \deleted{\unsure{que}} par le théorème de conjugaison que $\mathcal{C}$ s'identifie à $G/N$, et le morphisme $\psi:X\longrightarrow G$ n'est autre que celui défini dans le N° 2. On conclut alors par 2.6. b). Enfin, \deleted{il est aussi} cela montre aussi que $U$ est le \uline{plus grand} ouvert \inserted{$V$} de $G$ tel que $\psi^{-1}(V)\longrightarrow V$ soit un isomorphisme, car \deleted{il suffit de le vérifier encore dans} pour vérifier $V\subset U$, il suffit de le vérifier fibre par fibre.
}

\uline{Corollaire 3.2.} Sous les conditions de 3.1., soit $g$ une section régulière de $G$, i.e. telle que pour tout $s\in S$, $g(s)$ soit un point régulier de $G_s$. Alors il existe un et \inserted{un} seul sous-groupe de Cartan $C$ de $G$ \deleted{contenant $g(S)$} tel que $g$ soit une section de $C$.

En effet, l'hypothèse sur $g$ signifie que $g$ est une section de $U$, et la conclusion qu'il existe une unique section de $X$ qui la relève, ce qui \deleted{\unsure{revient}} n'est qu'une autre façon d'exprimer que $\psi^{-1}(U)\longrightarrow U$ est un isomorphisme.

Notons maintenant que l'ouvert $\psi^{-1}(U)$ du sous-groupe de Cartan $X$ de $G_{\deleted{\unsure{X}}\mathcal{C}}$ n'est autre que l'ouvert de $X$ formé des points de $X$ qui sont réguliers dans $G_{\deleted{\unsure{X}}\mathcal{C}}$ (par quoi on entend: réguliers dans leur fibre). On
}
\margin{\begin{tikzpicture}\node[draw,ellipse] {indice $\mathcal{C}$};\end{tikzpicture}}
\margin{128}
